% FIGURE AND LISTING GALLERY -- temporary, stage 2 only (editor-owned). Every figure F1 to F18 in its planned float
% type with its planned caption and label (PLAN.md section 4.2), grouped by chapter, and every snippet of every
% example in code/ (PLAN.md section 5), so that the author sees all the evidence at once before the prose is written.
% Removed from main.tex and deleted at stage 4, when the chapters carry the floats themselves. The labels here are
% the chapters' labels, so this file and a chapter must never be built together.
\chapter{Figure and listing gallery (stage 2, temporary)}\label{ch:gallery}

Every plot below reads a file in \code{data/} that \code{tools/generate_data.cpp} wrote by running the library;
every listing is quoted from a compiled example in \code{code/} that runs as a test. The captions are the planned
ones of \code{PLAN.md}; the chapters will carry the floats, and this gallery goes. Numbers in the captions are
\code{\val} macros from \code{data/values.tex}.

\section{Chapter 1: Why numerical derivatives}

\begin{marginfigure}
\centering
% F1 fig:secant (Chapter 1, margin): the tangent at x and two secants, a diagram without data. The curve is
% exp(t) around x = 0 with a step of 0.7; any smooth curve would do, and the text does not name it. The forward
% secant joins the points at x and x + h; the central secant joins the points at x - h and x + h and is nearly
% parallel to the tangent: the one thing to notice, in accent. The lines end at staggered abscissas so that their
% labels do not collide.
\begin{tikzpicture}[x=17mm, y=10.5mm, font=\sffamily\scriptsize]
  \pgfmathsetmacro{\stepH}{0.7}
  \pgfmathsetmacro{\fMinus}{exp(-\stepH)}
  \pgfmathsetmacro{\fPlus}{exp(\stepH)}
  \pgfmathsetmacro{\slopeF}{(\fPlus - 1) / \stepH}              % slope of the forward secant
  \pgfmathsetmacro{\slopeC}{(\fPlus - \fMinus) / (2 * \stepH)}  % slope of the central secant
  % the axis with the three nodes
  \draw[thin, black!55] (-1.2, -0.3) -- (1.25, -0.3);
  \foreach \nodeX/\nodeLabel in {-\stepH/{$x - h$}, 0/{$x$}, \stepH/{$x + h$}} {
    \draw[thin, black!55] (\nodeX, -0.3) -- ++(0, -0.05)
      node[below, black!70, inner sep=1.5pt] {\nodeLabel};
    \draw[densely dotted, black!40] (\nodeX, -0.3) -- (\nodeX, {exp(\nodeX)});
  }
  % the function
  \draw[line width=0.8pt, black!80, domain=-1.2:1.22, samples=60, smooth] plot (\x, {exp(\x)})
    node[left, black!80, inner sep=2pt] {$f$};
  % the tangent at x: slope f'(x) = 1
  \draw[line width=0.6pt, black!80] (-1.15, {1 - 1.15}) -- (0.95, {1 + 0.95})
    node[right, inner sep=1.5pt] {tangent};
  % the central secant through x - h and x + h
  \draw[line width=0.6pt, accent] (-1.15, {\fMinus + (-1.15 + \stepH) * \slopeC})
    -- (1.05, {\fMinus + (1.05 + \stepH) * \slopeC})
    node[right, inner sep=1.5pt] {central};
  % the forward secant through x and x + h
  \draw[line width=0.6pt, muted] (-0.3, {1 - 0.3 * \slopeF}) -- (1.15, {1 + 1.15 * \slopeF})
    node[right, inner sep=1.5pt] {forward};
  % the sampled points
  \fill[muted] (\stepH, \fPlus) circle (1.1pt);
  \fill[accent] (-\stepH, \fMinus) circle (1.1pt);
  \fill[black!80] (0, 1) circle (1.1pt);
\end{tikzpicture}

\caption{F1. A smooth curve with the tangent at $x$ and the secants through $(x, x+h)$ and $(x-h, x+h)$: the
central secant is nearly parallel to the tangent. A diagram without data.\label{fig:secant}}
\end{marginfigure}

The thread example of the book, nitrogen's enthalpy $\enth(T)$ in the Shomate form, is defined once in
\code{code/shomate.hpp} and quoted from there. Its heat capacity at $T_0 = \SI{1000}{K}$ is
$\cp = \SI{\valCpExact}{\joule\per\mole\per\kelvin}$ (the closed form, the answer key); \code{Central<2>} with the
default step of \SI{\valCpDefaultStep}{K} is \num{\valCpDefaultRelError} off relative, Ridders
\num{\valCpRiddersRelError} with an estimate of \SI{\valCpRiddersEstimate}{\joule\per\mole\per\kelvin} after
\valCpRiddersEvals{} evaluations.

\clearpage
\section{Chapter 2: Finite differences from Taylor's theorem}

\begin{figure}
\sidecaption{F2. Relative error of \code{Forward<1>} and \code{Central<2>} against the step $h$ on the thread,
$\enth(T)$ at $T_0$, for large steps: on logarithmic axes the errors fall on lines of slope 1 and 2, so halving the
step halves one error and quarters the other. Below \SI{0.01}{K} the central difference stops improving; why is the
subject of Chapter 3. Data \code{data/step-error-cp.csv}.\label{fig:fd-slopes}}
\centering
% F2 fig:fd-slopes (Chapter 2, text width): relative error of the forward difference (Forward<1>) and the central
% difference (Central<2>) against the step on the thread, H(T) at T0 = 1000 K, for the large steps 1e2 to 1e-2 K
% (DS1, data/step-error-cp.csv). On log axes the errors fall on lines of slope 1 and 2 while truncation dominates.
% The range stops at 1e-2 K, where the central difference stops improving (rounding, Chapter 3): the caption points
% ahead (editor, stage 2, decision S2-14). A domain restriction on a log axis is in natural logarithms:
% ln(1e-2) = -4.61, ln(1e2) = 4.61.
\begin{tikzpicture}
\begin{axis}[eqtufte, width=95mm, height=45mm, xmode=log, ymode=log,
    xlabel={step $h$ (K)}, ylabel={relative error},
    xmin=1e-2, xmax=1e2, ymin=1e-12, ymax=1e-1,
    xtickten={-2,-1,0,1,2}, ytickten={-12,-10,-8,-6,-4,-2},
    restrict x to domain=-4.62:4.62]
  \addplot[muted, mark=*] table[col sep=comma, x=h, y=forward1, comment chars={\#}] {data/step-error-cp.csv};
  \addplot[accent, mark=*] table[col sep=comma, x=h, y=central2, comment chars={\#}] {data/step-error-cp.csv};
  \node[muted, font=\sffamily\scriptsize, anchor=south east, align=right] at (axis cs:3, 1e-3)
    {forward difference: slope 1,\\ halving $h$ halves the error};
  \node[accent, font=\sffamily\scriptsize, anchor=north west, align=left] at (axis cs:0.3, 1e-9)
    {central difference: slope 2,\\ halving $h$ quarters the error};
\end{axis}
\end{tikzpicture}

\end{figure}

\clearpage
\section{Chapter 3: Rounding and the choice of step}

\begin{figure}
\sidecaption{F3. The V: the relative error of \code{Central<2>} against $h$ over the whole range of steps, the model
$E(h)$ dashed and the default step marked. The two branches, and the default step near the bottom. Data
\code{data/step-error-cp.csv}.\label{fig:v-curve}}
\centering
% F3 fig:v-curve (Chapter 3, text width): the V of the central difference (Central<2>) on the nitrogen example, H(T) at
% T0 = 1000 K: the measured relative error over the whole range of steps and the model
% E(h) = (h^2 |H'''| / 6 + eps |H| / h) / cp, dashed (DS1, data/step-error-cp.csv); the default step marked from
% \valCpDefaultStep and \valCpDefaultRelError (DS1b). The model exceeds 1 at the smallest steps, cut at the top.
\begin{tikzpicture}
\begin{axis}[eqtufte, width=95mm, height=45mm, xmode=log, ymode=log,
    xlabel={step $h$ (K)}, ylabel={relative error},
    xmin=1e-13, xmax=1e2, ymin=1e-12, ymax=1,
    xtickten={-12,-10,-8,-6,-4,-2,0,2}, ytickten={-12,-9,-6,-3,0},
    restrict y to domain=-30:0.01, unbounded coords=jump]
  \addplot[muted, densely dashed, line width=0.6pt] table[col sep=comma, x=h, y=model_central2,
    comment chars={\#}] {data/step-error-cp.csv};
  \addplot[accent, mark=*] table[col sep=comma, x=h, y=central2, comment chars={\#}] {data/step-error-cp.csv};
  % the default step: a thin muted line over the height of the axis, labelled at the top
  \draw[muted, thin] ({axis cs:\valCpDefaultStep,1}|-{rel axis cs:0,0})
    -- ({axis cs:\valCpDefaultStep,1}|-{rel axis cs:0,1})
    node[above, font=\sffamily\scriptsize, inner sep=1.5pt] {default step};
  \fill[accent] (axis cs:\valCpDefaultStep,\valCpDefaultRelError) circle (1.4pt);
  % the two branches and the model
  \node[font=\sffamily\scriptsize, black!70, anchor=south west, align=left] at (axis cs:3e-13, 1e-10)
    {rounding:\\ error ${}\propto 1/h$};
  \node[font=\sffamily\scriptsize, black!70, anchor=north west, align=left] at (axis cs:0.3, 1e-9)
    {truncation:\\ error ${}\propto h^2$};
  \node[font=\sffamily\scriptsize, muted, anchor=south west] at (axis cs:1e-9, 1e-3) {model $E(h)$, dashed};
\end{axis}
\end{tikzpicture}

\end{figure}

\begin{marginfigure}
\centering
% F5 fig:float-double (Chapter 3, margin): the V of the central difference (Central<2>) on exp(x) at x = 1 in float
% and in double (DS3, data/step-error-float.csv). Each floor sits at about eps^(2/3) of its type, drawn as a dashed
% guide from \valModelCentralTwoFirst{Float,Double}. Steps 1 to 1e-8 (on a log axis the domain restriction is in
% natural logarithms: ln(1e-8) = -18.42).
\begin{tikzpicture}
\begin{axis}[eqtufte, width=36mm, height=32mm, xmode=log, ymode=log,
    xlabel={step $h$}, ylabel={relative error},
    xmin=1e-8, xmax=1, ymin=1e-12, ymax=1,
    xtickten={-8,-4,0}, ytickten={-12,-8,-4,0},
    unbounded coords=jump, restrict x to domain=-18.43:0.01]
  \draw[accent, densely dashed, line width=0.5pt] ({rel axis cs:0,0}|-{axis cs:1,\valModelCentralTwoFirstFloat})
    -- ({rel axis cs:1,0}|-{axis cs:1,\valModelCentralTwoFirstFloat})
    node[below left, font=\sffamily\scriptsize, inner sep=1pt] {$\varepsilon_{\mathrm{float}}^{2/3}$};
  \draw[accent, densely dashed, line width=0.5pt] ({rel axis cs:0,0}|-{axis cs:1,\valModelCentralTwoFirstDouble})
    -- ({rel axis cs:1,0}|-{axis cs:1,\valModelCentralTwoFirstDouble})
    node[above left, font=\sffamily\scriptsize, inner sep=1pt] {$\varepsilon_{\mathrm{double}}^{2/3}$};
  \addplot[muted, mark=*] table[col sep=comma, x=h, y=error_float, comment chars={\#}] {data/step-error-float.csv};
  \addplot[black!75, mark=*] table[col sep=comma, x=h, y=error_double, comment chars={\#}] {data/step-error-float.csv};
  \node[muted, font=\sffamily\scriptsize, anchor=south west] at (axis cs:5e-7, 0.12) {float};
  \node[black!75, font=\sffamily\scriptsize, anchor=south west] at (axis cs:1.2e-8, 1e-8) {double};
\end{axis}
\end{tikzpicture}

\caption{F5. The V of \code{Central<2>} for float and double on $\mathrm{e}^x$ at 1: the floor sits at about
$\epsmach^{2/3}$ of each type. Data \code{data/step-error-float.csv}.\label{fig:float-double}}
\end{marginfigure}

The margin figures of this chapter (F5 here; F4b and F17 on the next page) are placed by the layout engine where
the margin has room; in the chapter each sits beside the paragraph that discusses it.

\clearpage

\begin{figure}
\sidecaption{F4. The V for a narrow feature at a large $\abs{x}$: the default step $\epsmach^{1/3} \scale$ sits far
to the right of the minimum, and the step $\epsmach^{1/3} w$ built from the feature's width $w$ lands at it. The
default step assumes that $f$ varies on the scale $\max(\abs{x}, 1)$; replace $\scale$ by the scale on which $f$
varies. Data \code{data/step-scale.csv}.\label{fig:step-scale}}
\centering
% F4 fig:step-scale (Chapter 3, text width): the V of the central difference (Central<2>) for a narrow feature at a
% large |x|, tanh((x - x0) / w) with x0 = 1000 and w = 0.5 at x = x0 + 0.2 (DS4, data/step-scale.csv). The default
% step eps^(1/3) s lies far to the right of the minimum; the step eps^(1/3) w from the feature's width lands at it.
% Both steps and their errors come from \valScale{Default,Fit}{Step,RelError}.
\begin{tikzpicture}
\begin{axis}[eqtufte, width=95mm, height=45mm, xmode=log, ymode=log,
    xlabel={step $h$}, ylabel={relative error},
    xmin=1e-12, xmax=10, ymin=1e-12, ymax=1,
    xtickten={-12,-10,-8,-6,-4,-2,0}, ytickten={-12,-9,-6,-3,0}]
  \addplot[black!75, mark=*] table[col sep=comma, x=h, y=error, comment chars={\#}] {data/step-scale.csv};
  % the default step eps^(1/3) max(|x|, 1)
  \draw[muted, thin] ({axis cs:\valScaleDefaultStep,1}|-{rel axis cs:0,0})
    -- ({axis cs:\valScaleDefaultStep,1}|-{rel axis cs:0,1})
    node[above, font=\sffamily\scriptsize, inner sep=1.5pt] {default step $\varepsilon^{1/3} s$};
  \fill[muted] (axis cs:\valScaleDefaultStep,\valScaleDefaultRelError) circle (1.4pt);
  % the step from the scale w on which f varies
  \draw[accent, thin] ({axis cs:\valScaleFitStep,1}|-{rel axis cs:0,0})
    -- ({axis cs:\valScaleFitStep,1}|-{rel axis cs:0,1})
    node[above, font=\sffamily\scriptsize, inner sep=1.5pt] {$\varepsilon^{1/3} w$};
  \fill[accent] (axis cs:\valScaleFitStep,\valScaleFitRelError) circle (1.4pt);
\end{axis}
\end{tikzpicture}

\end{figure}

\begin{marginfigure}
\centering
% F4b fig:step-scale-small (Chapter 3, margin): the V of the central difference (Central<2>) for y ln y at
% y = 1e-6 (DS4b, data/step-scale-small.csv). Steps above y sample below zero, where the function is not finite, so
% the curve stops there (nan in the CSV) and the default step eps^(1/3) gets no value; the step eps^(1/3) y lands
% near the minimum. Steps, status and error from \valSmall{DefaultStep,DefaultStatus,FitStep,FitRelError}.
\begin{tikzpicture}
\begin{axis}[eqtufte, width=36mm, height=32mm, xmode=log, ymode=log,
    xlabel={step $h$}, ylabel={relative error},
    xmin=1e-20, xmax=1e-5, ymin=1e-12, ymax=1e-2,
    xtickten={-20,-15,-10,-5}, ytickten={-12,-8,-4},
    unbounded coords=jump]
  \addplot[black!75, mark=*] table[col sep=comma, x=h, y=error, comment chars={\#}] {data/step-scale-small.csv};
  \draw[muted, thin] ({axis cs:\valSmallDefaultStep,1}|-{rel axis cs:0,0})
    -- ({axis cs:\valSmallDefaultStep,1}|-{rel axis cs:0,1})
    node[anchor=south east, font=\sffamily\scriptsize, inner sep=1pt, align=right] {default step:\\ \valSmallDefaultStatus};
  \draw[accent, thin] ({axis cs:\valSmallFitStep,1}|-{rel axis cs:0,0})
    -- ({axis cs:\valSmallFitStep,1}|-{rel axis cs:0,1});
  \node[accent, font=\sffamily\scriptsize, anchor=east, inner sep=1.5pt]
    at ({axis cs:\valSmallFitStep,1}|-{rel axis cs:0,0.72}) {$\varepsilon^{1/3} y$};
  \fill[accent] (axis cs:\valSmallFitStep,\valSmallFitRelError) circle (1.2pt);
\end{axis}
\end{tikzpicture}

\caption{F4b. The same at a small $\abs{x}$ in the user's units ($y \ln y$ at $y = 10^{-6}$): the default step
exceeds $x$ and crosses zero, so the central rule fails, while $\epsmach^{1/3} y$ gives the expected accuracy. The 1 in
$\max(\abs{x}, 1)$ is in your units. Data \code{data/step-scale-small.csv}.\label{fig:step-scale-small}}
\end{marginfigure}

\begin{marginfigure}
\centering
% F17 fig:noise (Chapter 3, margin): the V of the central difference (Central<2>) on exp(x) at x = 1, clean and with
% the values perturbed by a relative noise eta = 1e-10 and 1e-8 (DS19, data/step-error-noise.csv, the perturbation
% of code/noise.hpp). The rounding branch lifts to eta |f| / h and the minimum moves right and up. The default step
% (\valRootEpsThreeDouble, since s = 1) and the step eta^(1/3) for eta = 1e-8 with its error (\valNoiseEightFit...)
% are marked. Steps 1 to 1e-8 (on a log axis the domain restriction is in natural logarithms: ln(1e-8) = -18.42).
\begin{tikzpicture}
\begin{axis}[eqtufte, width=36mm, height=32mm, xmode=log, ymode=log,
    xlabel={step $h$}, ylabel={relative error},
    xmin=1e-8, xmax=1, ymin=1e-12, ymax=1,
    xtickten={-8,-4,0}, ytickten={-12,-8,-4,0},
    restrict x to domain=-18.43:0.01]
  \addplot[muted, mark=*] table[col sep=comma, x=h, y=clean, comment chars={\#}] {data/step-error-noise.csv};
  \addplot[black!75, mark=*] table[col sep=comma, x=h, y=noise_ten, comment chars={\#}] {data/step-error-noise.csv};
  \addplot[accent, mark=*] table[col sep=comma, x=h, y=noise_eight, comment chars={\#}] {data/step-error-noise.csv};
  \draw[muted, thin] ({axis cs:\valRootEpsThreeDouble,1}|-{rel axis cs:0,0})
    -- ({axis cs:\valRootEpsThreeDouble,1}|-{rel axis cs:0,1})
    node[above, font=\sffamily\scriptsize, inner sep=1.5pt] {default};
  \fill[accent] (axis cs:\valNoiseEightFitStep,\valNoiseEightFitRelError) circle (1.2pt);
  \node[accent, font=\sffamily\scriptsize, anchor=south west, inner sep=1pt] at (axis cs:2e-5, 1e-3)
    {$\eta = 10^{-8}$};
  \node[black!75, font=\sffamily\scriptsize, anchor=north west, inner sep=1pt] at (axis cs:1e-7, 1e-6)
    {$10^{-10}$};
  \node[muted, font=\sffamily\scriptsize, anchor=south west, inner sep=1pt] at (axis cs:1e-7, 1e-9)
    {clean};
\end{axis}
\end{tikzpicture}

\caption{F17. The V of \code{Central<2>} on $\mathrm{e}^x$ at 1 with the function's values perturbed by a relative
noise $\noise = 10^{-10}$ and $10^{-8}$, against the clean curve: the rounding branch lifts to $\noise \abs{f} / h$ and
the minimum moves right and up. A solver converged to $10^{-8}$ gives about five digits, at a step a thousand times
the default. Data \code{data/step-error-noise.csv}.\label{fig:noise}}
\end{marginfigure}

\clearpage
\section{Chapter 4: Higher-order and one-sided stencils}

\begin{figure}
\sidecaption{F6. Relative error against $h$ for \code{Central<2>} to \code{Central<10>}, $\mathrm{e}^x$ at 1: steeper
truncation slopes, rounding branches that nearly coincide, and minima that slide up the common rounding line as the
best step grows. A higher order buys accuracy by letting the step grow; the minimum falls from about
$\epsmach^{2/3}$ to about $\epsmach^{10/11}$, three and a half decades. Data
\code{data/step-error-exp.csv}.\label{fig:order-ladder}}
\centering
% F6 fig:order-ladder (Chapter 4, text width). Relative error against the step for Central<2> to Central<10>,
% exp(x) at x = 1 (DS2, data/step-error-exp.csv). The truncation branches steepen with P; the rounding branches
% nearly coincide (S from 1 to 2.3), and the minima slide up the common rounding line eps / h (muted, dashed),
% because a higher order lets the step grow. The five curves have five colours, ordered dark to light with P so
% that they also read in greyscale, and each label takes its curve's colour (the author's release decision: the
% grey ladder was hard to read). The float and caption stay in the chapter.
\definecolor{ladderTwo}{HTML}{2B2B2B}
\definecolor{ladderFour}{HTML}{0B5FAE}
\definecolor{ladderSix}{HTML}{1B8F8A}
\definecolor{ladderEight}{HTML}{6AA84F}
\definecolor{ladderTen}{HTML}{C79A1E}
\begin{tikzpicture}
\begin{axis}[eqtufte, width=95mm, height=45mm,
  xmode=log, ymode=log, log basis x=10, log basis y=10,
  xmin=1e-13, xmax=1, ymin=1e-16, ymax=1,
  xtick={1e-12,1e-10,1e-8,1e-6,1e-4,1e-2,1}, ytick={1e-16,1e-12,1e-8,1e-4,1},
  xlabel={step $h$}, ylabel={relative error}, unbounded coords=jump]
  % the common rounding line eps / h (relative, since f = f' for exp): a guide, not data
  \addplot[black!45, dashed, thin, domain=1e-13:1e-3, samples=2] {2^(-52) / x};
  \node[black!55, font=\scriptsize, anchor=south west, inner sep=1pt] at (axis cs:3e-9, 1.5e-7) {$\epsmach / h$};
  \addplot[ladderTwo] table[col sep=comma, comment chars={\#}, x=h, y=central2] {data/step-error-exp.csv};
  \addplot[ladderFour] table[col sep=comma, comment chars={\#}, x=h, y=central4] {data/step-error-exp.csv};
  \addplot[ladderSix] table[col sep=comma, comment chars={\#}, x=h, y=central6] {data/step-error-exp.csv};
  \addplot[ladderEight] table[col sep=comma, comment chars={\#}, x=h, y=central8] {data/step-error-exp.csv};
  \addplot[ladderTen] table[col sep=comma, comment chars={\#}, x=h, y=central10] {data/step-error-exp.csv};
  % labels on the truncation branches, left of each curve in its colour (the curve is lower to the left)
  \node[ladderTwo, font=\scriptsize, anchor=east, inner sep=1pt] at (axis cs:0.17, 5e-3) {\code{Central<2>}};
  \node[ladderFour, font=\scriptsize, anchor=east, inner sep=1pt] at (axis cs:0.17, 3e-5) {\code{Central<4>}};
  \node[ladderSix, font=\scriptsize, anchor=east, inner sep=1pt] at (axis cs:0.17, 2e-7) {\code{Central<6>}};
  \node[ladderEight, font=\scriptsize, anchor=east, inner sep=1pt] at (axis cs:0.17, 1.6e-9) {\code{Central<8>}};
  \node[ladderTen, font=\scriptsize, anchor=east, inner sep=1pt] at (axis cs:0.17, 1.2e-11) {\code{Central<10>}};
\end{axis}
\end{tikzpicture}

\end{figure}

\begin{marginfigure}
\centering
% F7 fig:one-sided (Chapter 4, margin). Forward<2> against Central<2> and Forward<4> against Central<4>, exp(x) at
% x = 1 (DS2): the same truncation slopes, higher rounding floors for the one-sided rules (larger S). One-sided rules
% in accent, central rules muted.
\begin{tikzpicture}
\begin{axis}[eqtufte, width=35mm, height=34mm,
  xmode=log, ymode=log, log basis x=10, log basis y=10,
  xmin=1e-13, xmax=1, ymin=1e-16, ymax=10,
  xtick={1e-12,1e-8,1e-4,1}, ytick={1e-16,1e-12,1e-8,1e-4,1},
  xlabel={step $h$}, ylabel={relative error}, unbounded coords=jump]
  \addplot[muted] table[col sep=comma, comment chars={\#}, x=h, y=central2] {data/step-error-exp.csv};
  \addplot[muted] table[col sep=comma, comment chars={\#}, x=h, y=central4] {data/step-error-exp.csv};
  \addplot[accent] table[col sep=comma, comment chars={\#}, x=h, y=forward2] {data/step-error-exp.csv};
  \addplot[accent] table[col sep=comma, comment chars={\#}, x=h, y=forward4] {data/step-error-exp.csv};
  % In the margin width the two rules of one order lie on top of each other along the truncation branch, so each
  % pair carries one order label, and the family is given by the colour (the key in the free top-left corner).
  \node[accent, font=\tiny, anchor=north west, inner sep=1pt] at (axis cs:1.5e-13, 5) {\code{Forward<P>}};
  \node[muted, font=\tiny, anchor=north west, inner sep=1pt] at (axis cs:1.5e-13, 0.15) {\code{Central<P>}};
  \node[font=\tiny, anchor=east, inner sep=1pt] at (axis cs:0.004, 1e-5) {$\aorder = 2$};
  \node[font=\tiny, anchor=east, inner sep=1pt] at (axis cs:0.004, 3e-10) {$\aorder = 4$};
\end{axis}
\end{tikzpicture}

\caption{F7. \code{Forward<2>} against \code{Central<2>} and \code{Forward<4>} against \code{Central<4>}: the same
slopes, higher floors (a larger $\ampl$). Data \code{data/step-error-exp.csv}.\label{fig:one-sided}}
\end{marginfigure}

\begin{marginfigure}
\centering
% F8 fig:default-step-accuracy (Chapter 4, margin). Relative error at the default step against the accuracy order P
% for Central<2..10> and Forward<1..6>, exp(x), N = 1 (DS7, data/default-step-accuracy.csv): the median over eleven
% points in [0.5, 1], the whiskers the smallest and largest value. Diminishing returns beyond P = 4 or 6 in double.
% The rows with nan (a P the family lacks) are discarded, so the lines join the existing orders.
\begin{tikzpicture}
\begin{axis}[eqtufte, width=38mm, height=34mm,
  ymode=log, log basis y=10, xmin=0.5, xmax=10.5, ymin=1e-16, ymax=1e-7,
  xtick={1,2,...,10}, ytick={1e-16,1e-14,1e-12,1e-10,1e-8},
  xlabel={accuracy order $\aorder$}, ylabel={relative error}, unbounded coords=discard,
  error bars/y dir=both, error bars/y explicit, error bars/error bar style={thin}, error bars/error mark options={rotate=90, mark size=1pt}]
  \addplot[accent, mark=*, mark size=1pt, error bars/error bar style={accent!60}]
    table[col sep=comma, comment chars={\#}, x=P, y=central,
      y error plus expr=\thisrow{central_max}-\thisrow{central}, y error minus expr=\thisrow{central}-\thisrow{central_min}]
    {data/default-step-accuracy.csv};
  \addplot[muted, mark=square*, mark size=1pt, error bars/error bar style={black!35}]
    table[col sep=comma, comment chars={\#}, x=P, y=forward,
      y error plus expr=\thisrow{forward_max}-\thisrow{forward}, y error minus expr=\thisrow{forward}-\thisrow{forward_min}]
    {data/default-step-accuracy.csv};
  \node[accent, font=\tiny, anchor=north, inner sep=1pt] at (axis cs:8, 1e-15) {\code{Central<P>}};
  \node[muted, font=\tiny, anchor=south west, inner sep=1pt] at (axis cs:3.3, 3e-12) {\code{Forward<P>}};
\end{axis}
\end{tikzpicture}

\caption{F8. Relative error at the default step against $\aorder$ for \code{Central<P>} and \code{Forward<P>},
$\mathrm{e}^x$: the median over eleven points in $[0.5, 1]$, the whiskers the smallest and largest value.
Diminishing returns beyond 4 or 6 for double. Data
\code{data/default-step-accuracy.csv}.\label{fig:default-step-accuracy}}
\end{marginfigure}

\begin{figure}
\sidecaption{F9. The nodes of \code{Central<4>}, \code{Forward<2>} and \code{Backward<2>} for $\dorder = 1$ and
$\dorder = 2$ on a line: node 0 filled where the rule uses it. A diagram without data.\label{fig:stencil-nodes}}
\centering
% F9 fig:stencil-nodes (Chapter 4, text width). The nodes of Central<4>, Forward<2> and Backward<2> for N = 1 and
% N = 2 on a line (a diagram, no data): a filled dot is a node the rule samples, in units of h from x. Node 0 (the
% point x itself) is in accent, filled where the rule uses it and hollow where it does not. From deriv.hpp:
% Central<P> uses the P + N - 1 integers around 0, without 0 for odd N and with 0 for even N; Forward<P> uses
% 0 ... P + N - 1; Backward<P> is Forward<P> mirrored.
\begin{tikzpicture}[x=7.5mm, y=5.5mm, font=\footnotesize,
  used/.style={fill=black!75, draw=none, circle, inner sep=0pt, minimum size=4pt},
  centre used/.style={fill=accent, draw=none, circle, inner sep=0pt, minimum size=4pt},
  centre unused/.style={fill=white, draw=accent, thin, circle, inner sep=0pt, minimum size=4pt}]
  % one row per rule: label, the line with ticks, the sampled nodes; the centre node drawn last
  \foreach \row/\name/\order/\nodes/\usesCentre in {
      0/Central<4>/1/{-2,-1,1,2}/0,
      1/Central<4>/2/{-2,-1,0,1,2}/1,
      2/Forward<2>/1/{0,1,2}/1,
      3/Forward<2>/2/{0,1,2,3}/1,
      4/Backward<2>/1/{-2,-1,0}/1,
      5/Backward<2>/2/{-3,-2,-1,0}/1} {
    \draw[black!55, thin] (-3.6, -\row) -- (3.6, -\row);
    \foreach \k in {-3,...,3} \draw[black!55, thin] (\k, -\row) ++(0, -0.12) -- ++(0, 0.24);
    \node[anchor=east] at (-4.0, -\row) {\texttt{\name}, $\dorder = \order$};
    \foreach \k in \nodes { \ifnum\k=0 \else \node[used] at (\k, -\row) {}; \fi }
    \ifnum\usesCentre=1 \node[centre used] at (0, -\row) {}; \else \node[centre unused] at (0, -\row) {}; \fi
  }
  % node labels under the last row
  \foreach \k in {-3,...,3} \node[anchor=north, font=\scriptsize, black!70] at (\k, -5.3) {$\k$};
  \node[anchor=north, font=\scriptsize, black!70] at (0, -5.95) {node $n_k$: the sample is at $x + n_k h$};
  % the one thing to notice: the centre is not sampled by Central<4> for f'
  \node[accent, font=\scriptsize, anchor=west] at (3.8, 0) {$x$ not sampled};
  \node[accent, font=\scriptsize, anchor=west] at (3.8, -1) {$x$ sampled};
\end{tikzpicture}

\end{figure}

\clearpage
\section{Chapter 5: Second to fourth derivatives}

\begin{figure}
\sidecaption{F10. The V of \code{Central<2>} for $\dorder = 1$ to 4, $\mathrm{e}^x$ at 1: the rounding branch
steepens as $1/h^{\dorder}$ and the minimum rises. Why the fourth derivative gives four digits where the first gives
eleven. Data \code{data/step-error-orders.csv}.\label{fig:v-curve-orders}}
\centering
% F10 fig:v-curve-orders (Chapter 5, text width). The V of Central<2> for N = 1 to 4, exp(x) at x = 1 (DS5,
% data/step-error-orders.csv). The truncation branches coincide (every derivative of exp is exp, and the constants
% 1/6, 1/12, 1/4, 1/6 differ little); the rounding branch steepens as 1/h^N and the minimum rises with N: the fourth
% derivative (accent) gives about four digits where the first gives eleven. Errors above 1 are cut off.
\begin{tikzpicture}
\begin{axis}[eqtufte, width=95mm, height=45mm,
  xmode=log, ymode=log, log basis x=10, log basis y=10,
  xmin=1e-13, xmax=1, ymin=1e-13, ymax=1,
  xtick={1e-12,1e-10,1e-8,1e-6,1e-4,1e-2,1}, ytick={1e-12,1e-9,1e-6,1e-3,1},
  xlabel={step $h$}, ylabel={relative error}, restrict y to domain=-14:-0.05, unbounded coords=jump]
  \addplot[black!80] table[col sep=comma, comment chars={\#}, x=h, y=n1] {data/step-error-orders.csv};
  \addplot[black!65] table[col sep=comma, comment chars={\#}, x=h, y=n2] {data/step-error-orders.csv};
  \addplot[black!50] table[col sep=comma, comment chars={\#}, x=h, y=n3] {data/step-error-orders.csv};
  \addplot[accent] table[col sep=comma, comment chars={\#}, x=h, y=n4] {data/step-error-orders.csv};
  % labels below and to the right of the rounding branches, where nothing else is drawn
  \node[font=\scriptsize, anchor=north west, inner sep=1pt] at (axis cs:3e-10, 1e-9) {$\dorder = 1$};
  \node[font=\scriptsize, anchor=north west, inner sep=1pt] at (axis cs:3e-7, 1e-5) {$\dorder = 2$};
  \node[font=\scriptsize, anchor=east, inner sep=1pt] at (axis cs:1.2e-5, 3e-2) {$\dorder = 3$};
  \node[accent, font=\scriptsize, anchor=south west, inner sep=1pt] at (axis cs:3e-4, 8e-2) {$\dorder = 4$};
\end{axis}
\end{tikzpicture}

\end{figure}

\begin{figure}
\sidecaption{F11. Relative error at the default step against $\dorder$ for \code{Central<2>}, \code{Central<4>} and
Ridders, double and float, $\mathrm{e}^x$ (medians over eleven points); the model $\epsmach^{\aorder/(\aorder+\dorder)}$
as a guide line; the float $f''''$ point marked as the warning case. Data
\code{data/order-accuracy.csv}.\label{fig:order-accuracy}}
\centering
% F11 fig:order-accuracy (Chapter 5, text width). Relative error at the default step against the derivative order N
% for Central<2>, Central<4> and Ridders, in double (solid) and float (dashed), exp(x), medians over eleven points
% in [0.5, 1] (DS6, data/order-accuracy.csv). The dotted guides are the model eps^(2/(2+N)) of Central<2> for each
% type (2^-52 and 2^-23: exact constants). The float fourth derivatives are boxed in accent: the warning case.
\begin{tikzpicture}
\begin{axis}[eqtufte, width=76mm, height=48mm,
  ymode=log, log basis y=10, xmin=0.7, xmax=4.3, ymin=1e-15, ymax=1,
  xtick={1,2,3,4}, ytick={1e-15,1e-12,1e-9,1e-6,1e-3,1},
  xlabel={derivative order $\dorder$}, ylabel={relative error at the default step}]
  % the model of Central<2>: eps^(P/(P+N)) with P = 2, for double and float (dotted; named in the caption)
  \addplot[muted, dotted, domain=1:4, samples=25] {(2^(-52))^(2/(2+x))};
  \addplot[muted, dotted, domain=1:4, samples=25] {(2^(-23))^(2/(2+x))};
  % double, solid
  \addplot[black!80, mark=*, mark size=1.2pt] table[col sep=comma, comment chars={\#}, x=N, y=c2_double] {data/order-accuracy.csv};
  \addplot[black!80, mark=square*, mark size=1.2pt] table[col sep=comma, comment chars={\#}, x=N, y=c4_double] {data/order-accuracy.csv};
  \addplot[black!80, mark=triangle*, mark size=1.5pt] table[col sep=comma, comment chars={\#}, x=N, y=rid_double] {data/order-accuracy.csv};
  % float, dashed, the same marks
  \addplot[black!60, dashed, mark=*, mark size=1.2pt, mark options={solid}] table[col sep=comma, comment chars={\#}, x=N, y=c2_float] {data/order-accuracy.csv};
  \addplot[black!60, dashed, mark=square*, mark size=1.2pt, mark options={solid}] table[col sep=comma, comment chars={\#}, x=N, y=c4_float] {data/order-accuracy.csv};
  \addplot[black!60, dashed, mark=triangle*, mark size=1.5pt, mark options={solid}] table[col sep=comma, comment chars={\#}, x=N, y=rid_float] {data/order-accuracy.csv};
  % labels at the right end of the double curves; the float curves carry the same marks
  \node[font=\scriptsize, anchor=west, inner sep=1pt] at (axis cs:4.07, 1.2e-5) {\code{Central<2>}, double};
  \node[font=\scriptsize, anchor=west, inner sep=1pt] at (axis cs:4.07, 4.9e-8) {\code{Central<4>}, double};
  \node[font=\scriptsize, anchor=west, inner sep=1pt] at (axis cs:4.07, 6.8e-10) {\code{Ridders}, double};
  \node[font=\scriptsize, anchor=south, inner sep=1pt] at (axis cs:2.5, 4e-3) {float, dashed (same marks)};
  % the warning case: every fourth derivative in float
  \draw[accent, thin] (axis cs:3.85, 3e-4) rectangle (axis cs:4.15, 4e-2);
  \node[accent, font=\scriptsize, anchor=south] at (axis cs:4, 4.5e-2) {float $f''''$: warning};
\end{axis}
\end{tikzpicture}

\end{figure}

\clearpage
\section{Chapter 6: Ridders' extrapolation}

\begin{table*}
\caption{F12. Ridders' tableau for $\cp$ at $T_0$ from $h_0 = T_0 / 32$: the \valTableauLevelsUsed{} levels the
method used, the entries $\Aij{i}{j}$ above and their relative errors with the rounding floor $r_i$ below; the entry
Ridders returned (level \valTableauBestLevel, column \valTableauBestColumn) in the accent colour. Generated bodies
\code{data/table-ridders-tableau.tex} and \code{data/table-ridders-tableau-errors.tex}.\label{tab:ridders-tableau}}
\centering
% F12 tab:ridders-tableau (Chapter 6). Ridders' tableau for cp at T0 from h_0 = T0 / 32: the levels the method
% used, the entries A_{i,j} (first table) and their relative errors with the rounding floor r_i (second table).
% Both bodies are generated (DS8, data/table-ridders-tableau.tex and data/table-ridders-tableau-errors.tex, seven
% columns: \valTableauLevelsUsed); the entry Ridders returned is flagged in accent by the generator. Needs the full
% width (table*): nine and ten columns of twelve-digit numbers do not fit the 114 mm text block.
% The bodies are read with the editor's \tablebody (macros.tex), the TeX primitive \input: LaTeX's \input{...} runs
% file hooks after the file, whose tokens start a new table row before \bottomrule and give "Misplaced \noalign".
\footnotesize\setlength{\tabcolsep}{3.5pt}
\begin{tabular}{@{}rr*{7}{r}@{}}
\toprule
 & & \multicolumn{7}{c}{entry $\Aij{i}{j}$ in J/(mol K)} \\
\cmidrule(lr){3-9}
level $i$ & $h_i$ (K) & $j = 0$ & $j = 1$ & $j = 2$ & $j = 3$ & $j = 4$ & $j = 5$ & $j = 6$ \\
\midrule
\tablebody{ridders-tableau}
\bottomrule
\end{tabular}
\par\medskip
\begin{tabular}{@{}rr*{7}{r}r@{}}
\toprule
 & & \multicolumn{7}{c}{relative error $\abs{\Aij{i}{j} - \cp} / \cp$} & floor \\
\cmidrule(lr){3-9}
level $i$ & $h_i$ (K) & $j = 0$ & $j = 1$ & $j = 2$ & $j = 3$ & $j = 4$ & $j = 5$ & $j = 6$ & $r_i / \cp$ \\
\midrule
\tablebody{ridders-tableau-errors}
\bottomrule
\end{tabular}

\end{table*}

\begin{figure}
\sidecaption{F12, companion diagram. The shape of the tableau: column 0 holds the base differences $D_i$ at the
halved steps; every other entry is built from its left neighbour and that neighbour's parent one level up. A diagram
without data.\label{fig:ridders-tableau-diagram}}
\centering
% F12 companion (Chapter 6, text width): the shape of Ridders' tableau as a diagram (no data). Column 0 holds the
% base differences D_i at the halved steps h_0 / 2^i; every other entry is built from its left neighbour and that
% neighbour's parent one level up, in Burden's difference form A_{i,j} = A_{i,j-1} + (A_{i,j-1} - A_{i-1,j-1}) /
% (4^j - 1). The one entry in accent shows the two parents it depends on; the diagonal entries (the latest entry
% of each level, against which the returned entry is checked) are in the darker ink.
\begin{tikzpicture}[x=14mm, y=7mm, font=\footnotesize,
  entry/.style={inner sep=1.5pt, text=black!55},
  diagonal/.style={inner sep=1.5pt, text=black!85},
  live/.style={inner sep=1.5pt, text=accent},
  dep/.style={-{Stealth[length=3pt]}, accent, thin}]
  % the entries: level i downwards, column j to the right; the diagonal darker
  \foreach \i in {0,...,4}
    \foreach \j in {0,...,\i} {
      \ifnum\i=\j
        \node[diagonal] (A\i\j) at (\j, -\i) {$\Aij{\i}{\j}$};
      \else
        \node[entry] (A\i\j) at (\j, -\i) {$\Aij{\i}{\j}$};
      \fi
    }
  % the one to notice and its two parents
  \node[live] (A32) at (2, -3) {$\Aij{3}{2}$};
  \node[live] (A31) at (1, -3) {$\Aij{3}{1}$};
  \node[live] (A21) at (1, -2) {$\Aij{2}{1}$};
  \draw[dep] (A31) -- (A32);
  \draw[dep] (A21) -- (A32);
  \node[accent, anchor=north, font=\scriptsize] at (2, -4.6)
    {$\Aij{3}{2} = \Aij{3}{1} + \dfrac{\Aij{3}{1} - \Aij{2}{1}}{4^2 - 1}$};
  % the steps of the levels
  \foreach \i/\divisor in {0/1, 1/2, 2/4, 3/8, 4/16}
    \node[anchor=east, text=black!70, font=\scriptsize] at (-0.4, -\i)
      {level $\i$, step \ifnum\i=0 $h_0$\else $h_0 / \divisor$\fi};
  % column headings
  \node[anchor=south, text=black!70, font=\scriptsize, align=center] at (0, 0.45) {column 0\\(the base $D_i$)};
  \foreach \j in {1,...,4}
    \node[anchor=south, text=black!70, font=\scriptsize] at (\j, 0.45) {column \j};
  % the diagonal
  \node[anchor=west, text=black!85, font=\scriptsize, align=left] at (4.35, -3.5) {the diagonal $\Aij{i}{i}$:\\the latest entry\\of each level};
\end{tikzpicture}

\end{figure}

\begin{figure}
\sidecaption{F13. Relative error of $\Aij{i}{j}$ against the level $i$ over 16 levels, one curve per column $j$, the
rounding floor $r_i$ dashed: column 0 descends slowly while each higher column reaches the floor within a few levels;
the floor doubles per level, so the entries then drift up. Ridders' stop and its returned entry are marked. The
extrapolation reaches in a few levels what the base difference would need a dozen for; after that, more levels only
add rounding. Data \code{data/ridders-tableau.csv}.\label{fig:ridders-convergence}}
\centering
% F13 fig:ridders-convergence (Chapter 6, text width). Relative error of the tableau entries A_{i,j} against the
% level i over 16 levels for cp at T0 (DS8, data/ridders-tableau.csv, long form: one row per (i, j); each curve
% selects its column j with a y expression and discards the other rows). Column 0 (the base difference) descends
% slowly; each higher column reaches the rounding floor r_i (dashed) within a few levels; the floor doubles per
% level, so the entries then drift up. Ridders' stop and its returned entry come from \valTableauStopLevel,
% \valTableauBestLevel and \valTableauBestRelError (the entry matched bit for bit by the generator).
\begin{tikzpicture}
\begin{axis}[eqtufte, width=95mm, height=48mm,
  ymode=log, log basis y=10, xmin=0, xmax=15.5, ymin=1e-16, ymax=1e-4,
  xtick={0,2,...,14}, ytick={1e-16,1e-14,1e-12,1e-10,1e-8,1e-6,1e-4},
  xlabel={level $i$ (step $h_0 / 2^i$, $h_0 = T_0/32$)}, ylabel={relative error of $\Aij{i}{j}$},
  unbounded coords=discard]
  % the rounding floor r_i (one value per level, taken from the j = 0 rows)
  \addplot[muted, dashed] table[col sep=comma, comment chars={\#}, x=i,
    y expr={\thisrow{j} == 0 ? \thisrow{floor} : nan}] {data/ridders-tableau.csv};
  \node[muted, font=\scriptsize, anchor=south, inner sep=1pt] at (axis cs:12.5, 3.5e-11) {floor $r_i$};
  % columns 0 to 4
  \addplot[black!85, mark=*, mark size=1pt] table[col sep=comma, comment chars={\#}, x=i,
    y expr={\thisrow{j} == 0 ? \thisrow{error} : nan}] {data/ridders-tableau.csv};
  \addplot[black!70, mark=square*, mark size=1pt] table[col sep=comma, comment chars={\#}, x=i,
    y expr={\thisrow{j} == 1 ? \thisrow{error} : nan}] {data/ridders-tableau.csv};
  \addplot[black!55, mark=triangle*, mark size=1.2pt] table[col sep=comma, comment chars={\#}, x=i,
    y expr={\thisrow{j} == 2 ? \thisrow{error} : nan}] {data/ridders-tableau.csv};
  \addplot[black!45, mark=diamond*, mark size=1.2pt] table[col sep=comma, comment chars={\#}, x=i,
    y expr={\thisrow{j} == 3 ? \thisrow{error} : nan}] {data/ridders-tableau.csv};
  \addplot[black!35, mark=pentagon*, mark size=1.2pt] table[col sep=comma, comment chars={\#}, x=i,
    y expr={\thisrow{j} == 4 ? \thisrow{error} : nan}] {data/ridders-tableau.csv};
  \node[font=\scriptsize, anchor=south west, inner sep=1pt] at (axis cs:7.3, 6e-9) {$j = 0$: the base difference $D_i$};
  \node[font=\scriptsize, anchor=south west, inner sep=1pt] at (axis cs:3.2, 3e-11) {$j = 1$};
  \node[font=\scriptsize, anchor=north west, inner sep=1pt] at (axis cs:6.3, 7e-16) {$j = 2, 3, 4$ (they coincide)};
  % Ridders' stop and the entry it returned
  \draw[accent, thin, dashed] (axis cs:\valTableauStopLevel, 1e-16) -- (axis cs:\valTableauStopLevel, 1e-4);
  \node[accent, font=\scriptsize, anchor=south west, inner sep=1pt] at (axis cs:\valTableauStopLevel, 3e-5) {Ridders stops};
  \node[accent, draw, circle, thin, inner sep=0pt, minimum size=6pt] at (axis cs:\valTableauBestLevel, \valTableauBestRelError) {};
  \node[accent, font=\scriptsize, anchor=south east, inner sep=1pt] at (axis cs:3.85, 1.1e-13) {returned};
\end{axis}
\end{tikzpicture}

\end{figure}

\clearpage

\begin{figure*}
\centering
% F14 fig:ridders-coverage (Chapter 6, full width, figure*). Ridders' error estimate against the actual error over
% the random study (DS9): N = 1 in dark grey, N = 2 muted; absolute values (data/README.md: relative ones blow up
% near a zero of f^(N)). The diagonal (accent) is estimate = actual error: a point above it is covered. An actual
% error of exactly 0 (an exact stop) cannot be shown on a log axis and is clipped to eps |f^(N)| with
% eps = 2^-52; the caption counts those cases (\valRiddersZeroErrorFirst, \valRiddersZeroErrorSecond).
% The plot reads data/ridders-coverage-plot-N{1,2}.csv, the two columns it needs with the clipping already applied
% by the generator: with the full ten-column files and an x expression per point, pgfplots used 2.6 of TeX's
% 5 million words of memory for this one figure and the book did not build (editor, stage 2, decision S2-17).
\begin{tikzpicture}
\begin{axis}[eqtufte, width=150mm, height=68mm,
  xmode=log, ymode=log, log basis x=10, log basis y=10,
  xmin=1e-19, xmax=1e-8, ymin=1e-18, ymax=1e-8,
  xtick={1e-18,1e-16,1e-14,1e-12,1e-10,1e-8}, ytick={1e-18,1e-16,1e-14,1e-12,1e-10,1e-8},
  xlabel={actual error $\abs{\text{value} - f^{(\dorder)}(x)}$}, ylabel={error estimate \code{abs_error}}]
  \addplot[only marks, mark=*, mark size=0.5pt, black!35] table[col sep=comma, comment chars={\#},
    x=actual, y=estimate] {data/ridders-coverage-plot-N2.csv};
  \addplot[only marks, mark=*, mark size=0.5pt, black!80] table[col sep=comma, comment chars={\#},
    x=actual, y=estimate] {data/ridders-coverage-plot-N1.csv};
  \addplot[accent, thin, domain=1e-18:1e-8, samples=2] {x};
  \node[accent, font=\scriptsize, anchor=north west, inner sep=2pt, rotate=24] at (axis cs:2e-17, 2e-17) {estimate $=$ actual error};
  \node[font=\scriptsize, anchor=west, inner sep=1pt] at (axis cs:1.5e-18, 1e-13) {$\dorder = 1$};
  \node[font=\scriptsize, anchor=west, inner sep=1pt, text=black!55] at (axis cs:1.5e-18, 1e-10) {$\dorder = 2$ (muted)};
  \node[font=\scriptsize, anchor=south east, inner sep=1pt, text=black!70] at (axis cs:1e-8, 1.3e-18) {covered: above the diagonal};
\end{axis}
\end{tikzpicture}

\caption{F14. Ridders' error estimate against the actual error over the random study ($\dorder = 1$ dark,
$\dorder = 2$ muted; nine smooth functions on $[0.2, 3]$, \valRiddersPointsPerFunction{} points each, no noise, no
singularity within $h_0$), the diagonal drawn: points above it are covered, \SI{\valRiddersCoverageFirst}{\percent} of
the \valRiddersCases{} cases for $\dorder = 1$ and \SI{\valRiddersCoverageSecond}{\percent} for $\dorder = 2$. Every
third case is drawn (\valRiddersPlotCases{} per order; all enter the statistics). An actual error of exactly 0 (an
exact stop; \valRiddersExactCasesFirst{} and \valRiddersExactCasesSecond{} cases) is clipped to
$\epsmach \abs{f^{(\dorder)}}$. Data \code{data/ridders-coverage-plot-N1.csv},
\code{-N2.csv}.\label{fig:ridders-coverage}}
\end{figure*}

\clearpage

\begin{marginfigure}
\centering
% F14b fig:ridders-cost (Chapter 6, margin). Histogram of the evaluations of fn per call over the random study
% (DS9, data/ridders-coverage-N1.csv and -N2.csv): N = 1 (dark, even counts: 2 per level) and N = 2 (muted, odd
% counts: 2 per level plus fn(x)). Unit bins centred on the integers, so the two series interleave. The tail of
% N = 2 at 19 and 21 is the slower convergence of several functions; the 11 restarts of log and sqrt below x = 0.25 are a small part of the last bar (R1-decision T7).
\begin{tikzpicture}
\begin{axis}[eqtufte, width=38mm, height=30mm,
  xmin=5.5, xmax=22.5, ymin=0, ymax=2100,
  xtick={6,10,14,18,22}, ytick={0,1000,2000}, scaled y ticks=false,
  xlabel={evaluations of \code{fn}}, ylabel={cases}]
  \addplot[hist={bins=17, data min=5.5, data max=22.5}, fill=black!30, draw=none]
    table[col sep=comma, comment chars={\#}, y=evaluations] {data/ridders-coverage-N2.csv};
  \addplot[hist={bins=17, data min=5.5, data max=22.5}, fill=black!75, draw=none]
    table[col sep=comma, comment chars={\#}, y=evaluations] {data/ridders-coverage-N1.csv};
  \node[font=\tiny, anchor=south, inner sep=1pt] at (axis cs:14.5, 1900) {$\dorder = 1$};
  \node[font=\tiny, anchor=south, inner sep=1pt, text=black!55] at (axis cs:17.5, 1580) {$\dorder = 2$};
\end{axis}
\end{tikzpicture}

\caption{F14b. Evaluations of \code{fn} per call over the random study: $\dorder = 1$ dark, $\dorder = 2$ muted; the
tail of $\dorder = 2$ are the restarts of $\log$ and $\sqrt{\cdot}$ below $x = 0.25$. Data
\code{data/ridders-coverage-N1.csv}, \code{-N2.csv}.\label{fig:ridders-cost}}
\end{marginfigure}

\begin{marginfigure}
\centering
% F15 fig:ridders-edge (Chapter 6, margin). Ridders for log x at x = 2^-k against k with the default initial step
% s/32 = 1/32 (DS10, data/ridders-edge.csv): above, the evaluations of fn (each halving of x costs one more restart,
% about two evaluations); below, the relative error, which stays near rounding until the level budget runs short.
% The three measured limits are marked from the generated values: the budget of 64 evaluations is first exhausted
% at k = \valRiddersEdgeCapExponent, ten correct digits last hold at k = \valRiddersEdgeTenDigitsExponent, and a
% value is last returned at k = \valRiddersEdgeLimitExponent; beyond it the result is non_finite (the nan rows break
% the curve). The last CSV row (x = 0.01, a non-integer k, E10) is dropped by the x expression.
\begin{tikzpicture}
\begin{groupplot}[group style={group size=1 by 2, vertical sep=3mm, xticklabels at=edge bottom},
  eqtufte, width=36mm, height=22mm, xmin=2, xmax=40, xtick={10,20,30,40}, unbounded coords=jump]
  \nextgroupplot[ymin=0, ymax=70, ytick={16,32,48,64}, ylabel={evaluations}]
  \addplot[black!75] table[col sep=comma, comment chars={\#},
    x expr={\thisrow{k} == floor(\thisrow{k}) ? \thisrow{k} : nan}, y=evaluations] {data/ridders-edge.csv};
  \draw[muted, thin] (axis cs:\valRiddersEdgeCapExponent, 0) -- (axis cs:\valRiddersEdgeCapExponent, 70);
  \draw[muted, thin] (axis cs:\valRiddersEdgeTenDigitsExponent, 0) -- (axis cs:\valRiddersEdgeTenDigitsExponent, 70);
  \draw[accent, thin] (axis cs:\valRiddersEdgeLimitExponent, 0) -- (axis cs:\valRiddersEdgeLimitExponent, 70);
  \node[muted, font=\tiny, anchor=south east, inner sep=1pt] at (axis cs:\valRiddersEdgeCapExponent, 70) {cap};
  \node[muted, font=\tiny, anchor=south, inner sep=1pt, yshift=1.1em] at (axis cs:\valRiddersEdgeTenDigitsExponent, 70) {10 digits};
  \node[accent, font=\tiny, anchor=south west, inner sep=1pt] at (axis cs:\valRiddersEdgeLimitExponent, 70) {last value};
  \nextgroupplot[ymode=log, log basis y=10, ymin=1e-16, ymax=1, ytick={1e-15,1e-10,1e-5,1},
    ylabel={relative error}, xlabel={$k$ ($x = 2^{-k}$)}]
  \addplot[black!75] table[col sep=comma, comment chars={\#},
    x expr={\thisrow{k} == floor(\thisrow{k}) ? \thisrow{k} : nan}, y=rel_error] {data/ridders-edge.csv};
  \draw[muted, thin] (axis cs:\valRiddersEdgeCapExponent, 1e-16) -- (axis cs:\valRiddersEdgeCapExponent, 1);
  \draw[muted, thin] (axis cs:\valRiddersEdgeTenDigitsExponent, 1e-16) -- (axis cs:\valRiddersEdgeTenDigitsExponent, 1);
  \draw[accent, thin] (axis cs:\valRiddersEdgeLimitExponent, 1e-16) -- (axis cs:\valRiddersEdgeLimitExponent, 1);
  \node[accent, font=\tiny, anchor=north west, inner sep=1pt] at (axis cs:\valRiddersEdgeLimitExponent, 1e-12) {non-finite};
\end{groupplot}
\end{tikzpicture}

\caption{F15. Evaluations and relative error of Ridders for $\log x$ at $x = 2^{-k}$ against $k$: restarts cost
about two evaluations per halving; the budget is exhausted from $k = \valRiddersEdgeCapExponent$, ten correct digits
hold to $k = \valRiddersEdgeTenDigitsExponent$, a value is returned to $k = \valRiddersEdgeLimitExponent$ and the
result is non-finite beyond. Data \code{data/ridders-edge.csv}.\label{fig:ridders-edge}}
\end{marginfigure}

The margin figures F14b and F15 stand on a page of their own here because a margin float cannot sit beside the
full-width F12 or F14; the chapter keeps them a page apart from those.

\clearpage
\section{Chapter 7: Using the library}

\begin{marginfigure}
\centering
% F16 fig:newton (Chapter 7, section 7.4): Newton's method for T with H(T) = H(1000 K) from 800 K, the slope by
% Central<2> with the default step (E12). The residual |H(T) - H*| in J/mol against the iteration: the number of
% correct digits doubles per step until the residual reaches the rounding of H. Data: DS11, data/newton.csv.
\begin{tikzpicture}
\begin{axis}[eqtufte, width=38mm, height=34mm, ymode=log,
  xlabel={iteration}, ylabel={residual $|H(T) - H^\ast|$ (J/mol)},
  xtick distance=1, ytick={1e-12,1e-8,1e-4,1,1e4},
  ymin=1e-13, ymax=1e5, xmin=-0.2, xmax=4.4]
  \addplot[accent, mark=*, mark size=1pt]
    table[col sep=comma, x=iteration, y expr={abs(\thisrow{residual})}, comment chars={\#}] {data/newton.csv}
    node[pos=0.5, anchor=south west, font=\scriptsize, accent, align=left] {quadratic:\protect\\digits double}
    node[pos=1, left=1pt, font=\scriptsize, muted] {rounding of $H$};
\end{axis}
\end{tikzpicture}

\caption{F16. Newton's residual against the iteration for $T$ with $\enth(T) = H^\ast$ and a \code{Central<2>}
slope: quadratic convergence until rounding, \valNewtonIterations{} iterations from \SI{\valNewtonStart}{K}. Data
\code{data/newton.csv}.\label{fig:newton}}
\end{marginfigure}

\begin{marginfigure}
\centering
% F18 fig:kink (Chapter 7, section 7.8): the published enthalpy of nitrogen across the range break at 500 K, where
% the two Shomate fits do not meet (H jumps by 0.77 J/mol). Relative error of cp at x = 500 K + delta against delta,
% default steps: Central<2>, Central<4>, Ridders' value and Ridders' relative error estimate (dashed). While the break
% lies inside a method's span the error is large and no method fails; only Ridders' estimate shows it. Data: DS20,
% data/kink.csv (E19). Labels sit on their curves (pos is a fraction of the drawn length), right of the drop of Central<2>.
\begin{tikzpicture}
\begin{axis}[eqtufte, width=38mm, height=56mm, xmode=log, ymode=log,
  xlabel={distance $\delta$ from the break (K)}, ylabel={relative error},
  xtick={1e-4,1e-2,1,1e2}, ytick={1e-14,1e-10,1e-6,1e-2,1e2},
  ymin=1e-15, ymax=1e2, xmin=4e-5, xmax=600, unbounded coords=jump]
  \addplot[muted]
    table[col sep=comma, x=delta, y=central2, comment chars={\#}] {data/kink.csv}
    node[pos=0, anchor=south west, font=\scriptsize] {\code{Central<2>}};
  \addplot[muted, densely dotted]
    table[col sep=comma, x=delta, y=central4, comment chars={\#}] {data/kink.csv}
    node[pos=0.15, anchor=south west, font=\scriptsize] {\code{Central<4>}};
  \addplot[black!80]
    table[col sep=comma, x=delta, y=ridders, comment chars={\#}] {data/kink.csv}
    node[pos=0.11, anchor=north west, font=\scriptsize] {Ridders};
  \addplot[accent, dashed]
    table[col sep=comma, x=delta, y=ridders_estimate, comment chars={\#}] {data/kink.csv}
    node[pos=0.11, anchor=south west, inner ysep=1pt, font=\scriptsize] {its estimate};
\end{axis}
\end{tikzpicture}

\caption{F18. Relative error of \code{Central<2>}, \code{Central<4>} and Ridders (value and estimate) against the
distance of $x$ from the range break at \SI{\valKinkJunction}{K}, where the two published Shomate fits of nitrogen
do not meet and $\enth$ jumps by \SI{\valKinkJumpEnthalpy}{\joule\per\mole}. Inside a method's span the error of a
straddling stencil grows as $1/h$ whatever $\aorder$ (for a kink it would be $\order(1)$ or $\order(h)$), the rule
returns it without failing, and only Ridders' estimate shows it. Evaluable is not smooth: the estimate flags it, the
stencil does not. Data \code{data/kink.csv}.\label{fig:kink}}
\end{marginfigure}

\clearpage
\section{Listings: every snippet of every example}

Each heading names the example of \code{PLAN.md} section 5 and its file; each listing is the snippet between the
markers named in the run-in heading.

\subsection*{Shared headers: \code{code/shomate.hpp}, \code{code/noise.hpp}}
\paragraph{shomate.} \snippetheader{shomate}{shomate}
\paragraph{shomate-H.} \snippetheader{shomate}{shomate-H}
\paragraph{shomate-cp.} \snippetheader{shomate}{shomate-cp}
\paragraph{noise.} \snippetheader{noise}{noise}

\subsection*{E1: \code{code/01-why-first.cpp}}
\paragraph{first.} \snippet{01-why-first}{first}
\paragraph{first-ridders.} \snippet{01-why-first}{first-ridders}

\subsection*{E2: \code{code/02-taylor-steps.cpp}}
\paragraph{forward-vs-central.} \snippet{02-taylor-steps}{forward-vs-central}

\subsection*{E3: \code{code/03-rounding-step.cpp}}
\paragraph{explicit-step.} \snippet{03-rounding-step}{explicit-step}
\paragraph{too-small.} \snippet{03-rounding-step}{too-small}

\subsection*{E4: \code{code/03-rounding-scale.cpp}}
\paragraph{narrow-feature.} \snippet{03-rounding-scale}{narrow-feature}
\paragraph{small-x.} \snippet{03-rounding-scale}{small-x}

\subsection*{E17: \code{code/03-rounding-noise.cpp}}
\paragraph{noisy-function.} \snippet{03-rounding-noise}{noisy-function}
\paragraph{noisy-coverage.} \snippet{03-rounding-noise}{noisy-coverage}

\subsection*{E5: \code{code/04-stencils-orders.cpp}}
\paragraph{higher-order.} \snippet{04-stencils-orders}{higher-order}

\subsection*{E6: \code{code/04-stencils-edge.cpp}}
\paragraph{validity-edge.} \snippet{04-stencils-edge}{validity-edge}

\subsection*{E7: \code{code/05-orders-second.cpp}}
\paragraph{second-derivative.} \snippet{05-orders-second}{second-derivative}
\paragraph{order-enum.} \snippet{05-orders-second}{order-enum}

\subsection*{E8: \code{code/05-orders-float.cpp}}
\paragraph{float-fourth.} \snippet{05-orders-float}{float-fourth}

\subsection*{E9: \code{code/06-ridders-basic.cpp}}
\paragraph{ridders-cp.} \snippet{06-ridders-basic}{ridders-cp}

\subsection*{E10: \code{code/06-ridders-edge.cpp}}
\paragraph{edge-restart.} \snippet{06-ridders-edge}{edge-restart}

\subsection*{E11: \code{code/07-using-expected.cpp}}
\paragraph{names.} \snippet{07-using-expected}{names}
\paragraph{expected-basics.} \snippet{07-using-expected}{expected-basics}
\paragraph{expected-chain.} \snippet{07-using-expected}{expected-chain}
\paragraph{estimate-pitfalls.} \snippet{07-using-expected}{estimate-pitfalls}

\subsection*{E12: \code{code/07-using-newton.cpp}}
\paragraph{newton.} \snippet{07-using-newton}{newton}
\paragraph{newton-call.} \snippet{07-using-newton}{newton-call}

\subsection*{E13: \code{code/07-using-sensitivity.cpp}}
\paragraph{rate.} \snippet{07-using-sensitivity}{rate}
\paragraph{sensitivity.} \snippet{07-using-sensitivity}{sensitivity}
\paragraph{scaled.} \snippet{07-using-sensitivity}{scaled}
\paragraph{jacobian.} \snippet{07-using-sensitivity}{jacobian}

\subsection*{E14: \code{code/07-using-multiprecision.cpp}}
\paragraph{multiprecision.} \snippet{07-using-multiprecision}{multiprecision}

\subsection*{E15: \code{code/07-using-constexpr.cpp}}
\paragraph{constexpr.} \snippet{07-using-constexpr}{constexpr}
\paragraph{constexpr-failure.} \snippet{07-using-constexpr}{constexpr-failure}
\paragraph{constexpr-value.} \snippet{07-using-constexpr}{constexpr-value}

\subsection*{E16: \code{code/08-api-usage.cpp}}
\paragraph{api-overloads.} \snippet{08-api-usage}{api-overloads}
\paragraph{api-rejected.} \snippet{08-api-usage}{api-rejected}

\subsection*{E18: \code{code/07-using-check.cpp}}
\paragraph{check-digits.} \snippet{07-using-check}{check-digits}
\paragraph{check-by-steps.} \snippet{07-using-check}{check-by-steps}
\paragraph{check-by-order.} \snippet{07-using-check}{check-by-order}
\paragraph{check-by-ridders.} \snippet{07-using-check}{check-by-ridders}
\paragraph{check-absolute.} \snippet{07-using-check}{check-absolute}

\subsection*{E19: \code{code/07-using-kink.cpp}}
\paragraph{kink.} \snippet{07-using-kink}{kink}
\paragraph{kink-outside.} \snippet{07-using-kink}{kink-outside}

\subsection*{E20: \code{code/07-using-solver.cpp}}
\paragraph{implicit-function.} \snippet{07-using-solver}{implicit-function}
\paragraph{implicit-derivative.} \snippet{07-using-solver}{implicit-derivative}
